\originalchapter{微正则系综与熵}\label{ch:micro}
\originalsection{相空间与等先验概率}
经典$N$粒子系统的一瞬间由所有坐标$q_i$和共轭动量$p_i$指定, 这些数值组成相空间中的一点. 相空间密度$\rho(p,q,t)$描述一族同样宏观准备的系统, 并非实空间的物质密度. Hamilton方程$\dot q_i=\partial H/\partial p_i$, $\dot p_i=-\partial H/\partial q_i$使相空间速度散度为零, 因为每对混合二阶导数相消. 连续性方程因此给出Liouville定理$\dd\rho/\dd t=0$. 若$\rho$只依赖守恒的$H$, 它就是一个静止的系综密度. 这不证明任一轨道遍历全部能面, 也不自动导出热平衡.

把相空间体积换成无量纲态数时需要量子作用量尺度. $h$是Planck（普朗克）常量, 单位为能量乘时间（SI中为$\mathrm{J\,s}$）, $\hbar=h/(2\pi)$是约化Planck常量. 它们表示同一量子尺度的两种约定: 频率$\nu$以每秒周期数计, 角频率$\omega=2\pi\nu$, 一个光子的能量为$h\nu=\hbar\omega$. 相空间中每对坐标与共轭动量用$h$归一, 故三维$N$粒子需要$h^{3N}$.

固定$E,V,N$的孤立系统, 把能量限制在宽度$\Delta$的壳层$E<H\le E+\Delta$. 微正则\textbf{假设}是可达微观状态等概率. 壳层包含的无量纲状态数写为
\[
\Omega(E,V,N)=\frac1{N!h^{3N}}\int_{E<H\le E+\Delta}\dd^{3N}p\dd^{3N}q.
\]
$h$使相空间体积无量纲, $N!$处理同种不可分辨粒子的经典计数; 固定在不同晶格位置的振子可区分, 则不除$N!$. 量子系统直接数满足能窗的能量本征态. 能窗必须包含许多微观能级, 又远小于宏观能量变化尺度.

定义累计状态数$\Phi(E)=\int_{H\le E}\dd\Gamma$和态密度$D(E)=\partial_E\Phi(E)$, 则窄窗内$\Omega\simeq D(E)\Delta$. 此处$\dd\Gamma$简记含上述归一因子的相空间积分元, 与第1章定义的Gamma函数$\Gamma(a)$不同. $D$每单位能量计数, 与$\Omega$不同. 对通常短程体系, 壳层与体积熵的差只有次广延修正, 在常规热力学极限中给出同样主阶结果; 小体系、受限能谱或负温体系不能无条件把它们等同. 本册温度采用壳层Boltzmann熵的定义.

若只观测部分变量$X$, 则$p(X)$等于与$X$相容的状态数除以总状态数. 从庞大集合中缩减自由度, 就把精确能量约束变成子系统的概率分布.

\originalsection{理想气体的状态计数}
单原子理想气体有$H=\sum_{i=1}^{3N}p_i^2/(2m)$, 位置积分给$V^N$. 动量约束是半径$\sqrt{2mE}$的$3N$维球, 于是
\[
\Phi=\frac{V^N(2\pi mE)^{3N/2}}{N!h^{3N}\Gamma(3N/2+1)}.
\]
球体积公式可由Gaussian积分说明: $\int_{\mathbb R^d}\e^{-r^2}\dd^dx=\pi^{d/2}$, 同时等于单位球面积乘$\Gamma(d/2)/2$, 再除$d$得单位球体积$\pi^{d/2}/\Gamma(d/2+1)$. 因此奇数维也能使用, 不把半整数阶乘当普通整数阶乘.

$S=\kb\ln\Omega$把可达状态的乘法变成可加的宏观量. 用Stirling近似保留主阶, 可得Sackur--Tetrode式
\[
S=N\kb\left[\ln\left\{\frac VN\left(\frac{4\pi mE}{3Nh^2}\right)^{3/2}\right\}+\frac52\right]+o(N).
\]
$E$和$V$、$N$一起放大时括号内保持不变, 显示广延性. 遗漏$N!$仍能在固定$N$时给出正确压强, 但会破坏混合熵和化学势, 不能因前几道题算对就认为计数正确.

设两系统弱耦合、总能量固定, 能量分配$E_1$的概率正比$\Omega_1(E_1)\Omega_2(E-E_1)$. 最概然值满足$\partial_{E_1}S_1=\partial_{E_2}S_2$, 所以定义$(\partial S/\partial E)_{V,N}=1/T$. 对气体得到$E=3N\kb T/2$. 由$dS=dE/T+P\dd V/T$得$(\partial S/\partial V)_{E,N}=P/T$, 从而$PV=N\kb T$.

选择某粒子的$x$动量, 剩余动量有能量$E-p_x^2/(2m)$. 壳层积分给$p(p_x)\propto(E-p_x^2/(2m))^{(3N-3)/2}$, 支持区为$|p_x|\le\sqrt{2mE}$. 在$N\gg1$且该分量携带能量远小于$E$时, 对数展开$\ln(1-u)=-u+O(u^2)$, 于是
\[
p(p_x)\simeq\frac1{\sqrt{2\pi m\kb T}}\e^{-p_x^2/(2m\kb T)}.
\]
有限系统的动量分量受总能量约束而相关, Gaussian独立性是热力学极限中小子系统的近似.

\originalsection{微正则振子到 Boltzmann 边缘}
固定各位置的 $N$ 个同频经典振子，能量 $\epsilon_i=p_i^2/(2m)+m\omega^2q_i^2/2$。把每对坐标缩放成圆形，$2N$ 维球的体积给无量纲累计态数
\[\Phi_N(E)=\frac1{N!}\left(\frac E{\hbar\omega}\right)^N,\qquad D_N(E)=\frac{E^{N-1}}{(N-1)!(\hbar\omega)^N}.\]
这里的 $N!$ 来自正能量单纯形或球体积，\textbf{不是}不可分辨粒子的置换校正；固定振子各有各的位置标签。在总能窄壳中，指定一振子能量 $e$ 后，余者可达态数 $D_{N-1}(E-e)$；因一振子的能量 DOS 为 $1/(\hbar\omega)$，$N>1$ 时
\[p(e)=\frac{D_{N-1}(E-e)}{\hbar\omega D_N(E)}=\frac{N-1}E\left(1-\frac eE\right)^{N-2},\quad0\le e\le E.\]
积分为一，均值 $E/N$。记 $\beta=1/(\kb T)$ 为热能尺度的倒数，单位为能量的倒数。取 $N\to\infty$、$E/N=\kb T$ 固定，$p(e)\to\beta e^{-\beta e}$。回到一振子的相空间，$p(p,q)\to\omega e^{-\beta\epsilon(p,q)}/(2\pi\kb T)$，两变量分别高斯。原8.044第12讲用这个具体例子把状态数比例接到下一章热库推导；完整有限尺寸式亦见作业四第6题。

\originalsection{第二定律与基本热力学关系}
平衡对应宏观约束下总熵的最大值. 大系统的状态数差随$N$指数放大, 所以解除约束后通常观察到走向更大熵的宏观状态. 这是关于宏观粗粒化和典型演化的统计定律; Liouville定理仍保持精细相空间密度的Gibbs熵, 有限孤立系统也可能有涨落与复现. 不把“末态最大熵”冒充每个微观轨道上严格单调增加的定理.

系统与固定温度热库交换热$Q$, 总熵变化$\Delta S_{\rm total}=\Delta S_{\rm sys}-Q/T_{\rm bath}\ge0$, 于是$\Delta S_{\rm sys}\ge Q/T_{\rm bath}$. 可逆平衡路径有$\delta Q_{\rm rev}=T\dd S$. 对单相平衡状态, 第一律和这个可逆路径关系结合为
\[
\dd U=T\dd S-P\dd V+\mu\dd N+\sum_j X_j\dd x_j.
\]
这是状态函数的关系, 用一条假想可逆路径就能计算不可逆过程的端点熵差; 并不宣称真实不可逆热量等于$T\dd S$. 化学势$\mu$的物理作用将在粒子交换时说明.

\originalsection{两能级与间隙原子}
$N$个可区分位置上的单元, 每个只有能量0和$\epsilon>0$两个状态. 激发数$N_1=E/\epsilon$固定时$\Omega=\binom N{N_1}$. Stirling式给
\[
S\simeq-N\kb[x\ln x+(1-x)\ln(1-x)],\quad x=N_1/N,
\quad\frac1T=\frac{\kb}{\epsilon}\ln\frac{1-x}{x}.
\]
故$x=1/(\e^{\epsilon/(\kb T)}+1)$, 单个单元激发概率也恰为$N_1/N$. 其热容为$N\kb(\epsilon/\kb T)^2\e^{\epsilon/\kb T}/(1+\e^{\epsilon/\kb T})^2$, 低温有能隙抑制, 高温因能量饱和而再降至零. 激发多于一半时壳层熵下降, 定义给负温; 这需要能谱有上界和隔离于无上界的自由运动, 并非比零开尔文更冷.

间隙原子模型把正常晶格位置和高能间隙位置都计入. 若$N$个原子占据$N$个正常位和$N_i$个间隙位, 有$n$个原子转移且每个代价$\epsilon$, 状态数为$\binom Nn\binom{N_i}n$. 两个因子分别选空位和已占间隙, 不能只数被移动的原子. 对数求导给
\[
\frac1T=\frac{\kb}{\epsilon}\ln\frac{(N-n)(N_i-n)}{n^2}.
\]
稀缺陷极限$n\ll N,N_i$下$n\simeq\sqrt{NN_i}\e^{-\epsilon/(2\kb T)}$. 指数中的二分之一来自一个激发同时产生空位与间隙的成对计数.
